\section{What must be known before a method can run}
\label{sec:method-requirement-language}

A company model can name a tax method while still lacking the information that
method needs. It may know the tax rate but not the relevant tax bases, or have a
reported consolidated balance without an allocation to the legal book being
simulated. Naming the implementation does not resolve either problem. The
method definition therefore needs to say which quantities and findings it
requires, how they are related, and what prevents the method from proceeding.

The preceding portable contract already calls for typed preconditions. The
\code{MethodRequirements/1} reference gives a small part of that contract a
common executable meaning. A requirement has an identity, a typed predicate,
the exact input identities on which it depends, references to required admission
findings, and different consequences for a false or unresolved predicate.
The same interpreter can assess declarations for several methods. It does not
replace their accounting algorithms, source checks or state transitions.

Consider the conditional annual-company example. Its floating-rate share is a
fraction of the declared debt principal, so the selected method requires a
value between zero and one. It also restricts the coupon interval to one year
and the opening facility draw and loan fee to zero. A negative fraction or a
share exceeding one contradicts the declared input domain. A nonzero fee instead
requires a method that handles fees. Those are different findings: the first
is invalid for the stated fraction; the second is outside the selected method's
scope. Neither can be repaired by silently changing the supplied number.

\subsection{A known contradiction survives an unrelated unknown}

Let $f$ denote the declared floating-rate share and $r$ the tax-rate input in
this finite annual profile. Its selected conditions include $0\leq f\leq1$
and $0\leq r\leq1$. These bounds belong to the declared profile; this example
does not assert a universal model of tax measurement. Suppose the caller gives
$f=2$ and leaves $r$ unknown. The supplied share already contradicts its meaning
as a fraction. Uncertainty about tax does not remove that contradiction.

The common assessment evaluates every declared requirement and retains each
finding. Its summary gives a known invalid input precedence over unsupported
scope, and unsupported scope precedence over unresolved requirements. Successful
selected requirements remain visible as well. This ordering summarizes the
results; it does not discard the reasons for an interrupted assessment.

The existing \code{CompanyOperationsDraft} validator follows a different order
in this mixed case: it returns on a globally unknown input before reaching the
fraction-domain test. It therefore reports unresolved for the same inputs.
The new reference reports invalid while retaining the unresolved tax condition.
The original validator is preserved, and this difference is recorded explicitly.
The checked agreement concerns selected single-condition comparisons, rather
than equivalence of every possible mixed input to the old validator.

Each input is separately marked known or unknown with a reason. An absent
binding is unresolved; it is never zero or false. A declaration that refers to
an undeclared input is structurally invalid. The literal text \emph{unknown} can
still be a known text value because the binding status is separate from its
value. Caller findings also have their own role. A true finding records what
the caller supplied; the evaluator does not establish that an accounting
judgment is correct merely by reading a Boolean.

The Boolean operators use the following finite tables. Here $U$ means an
unresolved supplied input or predicate, not a probability and not a set of
financial histories.
\begin{center}
\begin{tabular}{cc|cc}
$A$ & $B$ & $\operatorname{all}(A,B)$ & $\operatorname{any}(A,B)$ \\
\hline
true & true & true & true \\
true & false & false & true \\
true & $U$ & $U$ & true \\
false & true & false & true \\
false & false & false & false \\
false & $U$ & false & $U$ \\
$U$ & true & $U$ & true \\
$U$ & false & false & $U$ \\
$U$ & $U$ & $U$ & $U$
\end{tabular}
\end{center}
Negation exchanges true and false and leaves $U$ unresolved. Every operand's
dependency is retained even when another operand determines the result. Thus a
false conjunction may still report the input missing from its other branch.
Repeated references identify the same input; they introduce no independent
draws or weights.

Arithmetic and comparisons with an unknown operand remain unknown in this
fragment. In particular, an unresolved $x$ leaves both $x=x$ and $x-x$ unresolved
here. This conservative rule is deliberately different from symbolic deduction
over an explicitly inhabited domain or the coefficient cancellation used in the
source-relation calculation. Those procedures have additional premises. The
requirement interpreter evaluates supplied inputs and does not decide entailment
over admissible company histories.

\subsection{The same requirement in text and typed form}

Types carry explicit context maps. A rational amount can identify its currency,
scale, entity, book, measurement date and accounting basis. Numeric operations
and comparisons require exactly matching types and contexts. A reported
consolidated balance of 100 and a legal-book balance of 100 are consequently not
interchangeable merely because their numerals agree. An explicit allocation or
conversion belongs to a separately specified method. Exact rational arithmetic
also describes the calculation, not the precision or truth of a reported amount.

The first expression fragment contains references, typed constants, equality,
ordered comparison, addition, subtraction and the Boolean operators. It has no
arbitrary calls, imports, multiplication, quantifiers or collection traversal.
These restrictions make the selected semantics inspectable while leaving the
broader method definition open to later compatible extensions.

\noindent\begin{minipage}{\linewidth}
Text uses named, semicolon-terminated declarations and prefix expressions. For
example, a requirement that a declared share does not exceed one can be written
as follows. The context and input roles remain explicit.
\begin{lstlisting}
input "share" rational assumption
  {"unit":"ratio"} "case:declared-share";
require "share.upper" invalid unresolved ["share"] []
  (le (ref "share")
      (const rational {"unit":"ratio"} "1/1"));
bind "share" known "1/2";
\end{lstlisting}
\end{minipage}

The display wraps declarations for reading; the reference text grammar places
each complete declaration on one line. A complete document also contains the
method identity and version, capability and residual-prerequisite declarations,
an exact selected manifest digest and an end marker. The lower bound and other
conditions must be declared separately; the upper-bound example supplies none
implicitly. Rational values use reduced numerator/denominator strings. Boolean
values are true or false, not integers or quoted words.

The frontend lowers these declarations to the same records used by the typed
library and calls the same evaluator. Strict exchange decoding rejects duplicate
JSON keys, including nested context keys. The text parser rejects duplicate
declarations, unknown operators, malformed types, missing terminators and
trailing content. Quoted delimiters and Unicode remain part of the exact supplied
text values, without normalization. These checks establish a bounded authoring
correspondence; they do not establish that users prefer this syntax.

\subsection{Content identity leaves authority with the caller}

The caller selects a manifest by its canonical SHA256 digest. Changing a
predicate, context, source role or residual prerequisite changes that identity.
A stale selection is rejected, and an unavailable definition or required
capability is reported separately. The manifest loads no code. Transporting an
availability list does not establish that an implementation is installed, and
a content hash does not authenticate the sender. Registry ownership and
capability availability remain responsibilities of the trusted caller.

A successful result is called \code{requirements\_satisfied} within the scope
\code{selected\_requirements\_only}. It retains the unassessed prerequisites
and explicitly leaves method eligibility, execution, source admission and
accounting admission false. For the checked company bridges, those remaining
conditions include source binding, account and state controls, event order,
unselected allocation equations and the actual method transition. A nonempty
source reference or an empty admission-reference list cannot settle them.

The common interpreter has 35 checked selected-condition comparisons with the
unchanged company-operation and return-module functions, complemented by
independent truth tables, exact arithmetic and adversarial type/context cases.
The separate frontend has exact typed/text and result comparisons across five
examples, including constructed legal-book, reported-measurement and
insurance-cohort contexts. The latter demonstrate language behavior, not new
insurance or source-admission methods. The mixed unknown/contradiction case
above remains a documented difference.

These references make prerequisites reusable and inspectable across the tested
profiles. General method applicability, account-state conversion, source-to-latent
reconstruction and independent acceptance still require their own evidence.
The following section extends this reference to exact record binding,
dimensional arithmetic and guarded applicability. Complete method boundaries
still require more than a partial predicate list before execution can proceed.
